\chapter{Motif Structure in the Score Graph}\label{ch:motif}
The attribution results of the previous chapters leave the score graph with
a modest role. On both corpora its value is temporal smoothing over about
two beats, and the pitch term adds nothing measurable
(Chapter~\ref{ch:discussion}, Appendix~\ref{app:trackscale}). This chapter
asks whether a different kind of score structure, one that the time window
cannot see, carries expressive information: the repetition of a motif. A
performer who shapes a motif at its first statement is expected to shape
its returns in a related way, and a return that has not yet been heard
should then be predictable from the statements that have. The chapter adds
one edge family to the score graph from motif annotations, states what the
family does to a hidden note, and evaluates it on Beethoven sonata
movements under a masking scheme that hides whole occurrences. The
development studies and a pre-specified evaluation on held-out performances
follow the protocol of \S\ref{sec:hygiene}. The evaluation uses score
features only, without the music model, since the question is one score
structure against another on identical features.

\section{An edge family from motif occurrences}\label{sec:motif-family}
The BPS-motif annotations \cite{hsiao2023bpsmotif} mark, in all thirty-two
first movements of the Beethoven piano sonatas, the occurrences of each
movement's motifs: an occurrence $O_k$ is a set of notes with a span
$[b_k^{\mathrm{start}}, b_k^{\mathrm{end}}]$ in beats and a type $q(k)$, and
two occurrences of one type are returns of one motif. Returns differ in
length and in note count, so there is in general no note-to-note map
between them. The family therefore works in a coordinate that every
occurrence shares, the relative position of a note inside its occurrence,
\begin{equation}
  u_i \;=\; \frac{b_i - b_k^{\mathrm{start}}}{\,b_k^{\mathrm{end}} - b_k^{\mathrm{start}}\,}
  \;\in\; [0, 1], \qquad i \in O_k,
  \label{eq:motif-u}
\end{equation}
and joins notes of two different occurrences of the same type by their
closeness in that coordinate,
\begin{equation}
  \bigl(W_{\mathrm{motif}}\bigr)_{ij} \;=\;
  \max_{k \neq l,\; q(k) = q(l)}\;
  \mathbf 1[i \in O_k]\, \mathbf 1[j \in O_l]\;
  \exp\!\Bigl(-\frac{(u_i - u_j)^2}{2\,\ell_u^2}\Bigr),
  \qquad \bigl(W_{\mathrm{motif}}\bigr)_{ii} = 0 .
  \label{eq:motif-adj}
\end{equation}
The length scale $\ell_u$ is fixed at a tenth of the span. Returns of equal
length then receive one-to-one edges between corresponding notes, and
returns of unequal length receive a soft block between the notes near the
same relative position. Nothing is aligned and nothing inside the family is
learned. A note that belongs to no occurrence, or whose type occurs only
once in the excerpt, has no motif edge.

The family enters the model of Chapter~\ref{ch:model} in two forms. In the
\textbf{shared form} the adjacency \eqref{eq:adjacency} gains the family
with one weight,
\begin{equation}
  W \;=\; W_{\mathrm{time}} + w_{\mathrm m}\, W_{\mathrm{motif}},
  \qquad \Lg = D - W,
  \qquad K_G(s) = U \diag\bigl(g(\nu; s)\bigr) U^{\top},
  \label{eq:motif-shared}
\end{equation}
so that one graph kernel, built from the Laplacian of the combined
adjacency, is shared by the three channels through the coregionalization
matrix $B$ of \eqref{eq:gpfirst}, and $\log w_{\mathrm m}$ joins the
hyperparameters fitted by the per-piece evidence. The evidence can drive
$w_{\mathrm m}$ to zero, which recovers the time-only graph exactly. In the
\textbf{per-channel form} the family is its own additive term with one
scale per channel,
\begin{equation}
  K \;=\; B \otimes K_G(s) \;+\; \diag(c_{\mathrm m}) \otimes K_{\mathrm m}
  \;+\; \sum_f \diag(c_f) \otimes X_f X_f^{\top}
  \;+\; \diag(\varsigma^2) \otimes I,
  \qquad K_{\mathrm m} = U_{\mathrm m} \diag\bigl(g(\nu_{\mathrm m}; 1)\bigr) U_{\mathrm m}^{\top},
  \label{eq:motif-pc}
\end{equation}
where $(\nu_{\mathrm m}, U_{\mathrm m})$ is the eigendecomposition of the
Laplacian of $W_{\mathrm{motif}}$ on the motif notes, $K_{\mathrm m}$ is zero
on every other note, and $c_{\mathrm m} \in \R^{3}$ is learned by the
evidence like the feature scales $c_f$. Since $K_{\mathrm m}$ factorizes as
$X_{\mathrm m} X_{\mathrm m}^{\top}$ with
$X_{\mathrm m} = U_{\mathrm m} \diag(g^{1/2})$, the per-channel form is a
feature kernel of \eqref{eq:gpfirst} and needs no change to the model's
inference. The two forms differ in what they let the evidence decide: the
shared form ties the three channels to one motif weight through $B$, and
the per-channel form lets a channel switch the family off on its own.

\section{What a motif edge does to a hidden note}\label{sec:motif-conditional}
For one channel with the additive kernel $g(\nu; s) = 1/(1 + s\nu)$ the
prior is the Gaussian Markov random field with precision $I + s \Lg$, and
the conditional of a hidden note $h$ given all others is explicit,
\begin{equation}
  \mathbb E\bigl[y_h \mid y_{-h}\bigr] \;=\; \frac{s \sum_{j} W_{hj}\, y_j}{1 + s\, d_h},
  \qquad
  \operatorname{Var}\bigl[y_h \mid y_{-h}\bigr] \;=\; \frac{1}{1 + s\, d_h},
  \qquad d_h = \sum_j W_{hj} .
  \label{eq:motif-cond}
\end{equation}
Splitting the degree into its two families,
$d_h = d_h^{\mathrm{time}} + w_{\mathrm m} d_h^{\mathrm{motif}}$, writes the
conditional mean as a mixture of two neighbourhood means,
\begin{equation}
  \mathbb E\bigl[y_h \mid y_{-h}\bigr] \;=\;
  \frac{s\, d_h^{\mathrm{time}}}{1 + s d_h}\;\overline{y}^{\,\mathrm{time}}_h
  \;+\; \frac{s\, w_{\mathrm m}\, d_h^{\mathrm{motif}}}{1 + s d_h}\;\overline{y}^{\,\mathrm{motif}}_h ,
  \label{eq:motif-mix}
\end{equation}
the mean of the temporal neighbours and the mean of the corresponding notes
in the other occurrences, weighted by relative position. The full posterior
of \eqref{eq:posterior} mixes the same two sources through $B$, the feature
kernels and the noise. Two readable predictions follow. If performers shape
a motif in a related way at each return, the motif mean is closer to the
hidden value than the temporal mean, the evidence raises $w_{\mathrm m}$,
the error falls and the predictive band tightens by the extra degree. If
they do not, the motif mean is noise, the evidence drives $w_{\mathrm m}$ to
zero and every motif arm coincides with the time-only graph. The second
outcome is the family's own falsification route and needs no separate
test.

\section{Evaluation design}\label{sec:motif-design}
The data are the ASAP first movements with an annotated counterpart. The
annotation's note list is matched to the ASAP score by crotchet onset and
pitch after a global offset search over the opening notes, and a movement
is used when at least $85\%$ of its excerpt notes match. Twenty-five
annotated movements have an ASAP score. Two of them sit in the piano
held-out test set and are excluded, and two more fall below the matching
threshold, which leaves twenty-one movements, each as the first $400$
matched notes of one aligned performance, the thesis's excerpt convention.
A note belongs to an occurrence when its onset lies in the occurrence's
span and its annotated type is the occurrence's type. The arms are the
time-only Gaussian graph, which the attribution studies single out, and the
same graph with the motif family in one of its two forms, paired on
identical masks and data. Every arm carries the score features with a
per-piece bias column, the coregionalization matrix, the floored noise, the
exact per-piece evidence and the predictive-variance floor of
\eqref{eq:floor}.

Two masking schemes are used. The random scheme is the thesis protocol,
$40\%$ of notes hidden at random. The occurrence scheme hides whole
occurrences: for every motif type with at least two occurrences, each
occurrence is hidden with probability $0.4$, with at least one occurrence of
the type left observed and at least one hidden, and every note outside a
motif observed. The occurrence scheme is the application's case, a passage
not yet played whose earlier statements have been heard, and it is where the
family can act, since under the random scheme the two-beat time window
already interpolates inside an occurrence. Scores are computed on the hidden
motif notes whose type has another occurrence with an observed note, per
channel, as the Gaussian log score and the root mean squared error, paired
per piece with seeds averaged, with $95\%$ bootstrap intervals over pieces.
The hypotheses were fixed before each run. The first study predicted a
lower timing log score for the motif arm against the time-only graph on the
hidden motif notes under the occurrence scheme, with the interval excluding
zero, and a lower timing error. The second study predicted that the
per-channel form keeps that gain and removes the velocity miscalibration
of the shared form. The third study predicted that edges built from a
model's predicted occurrences keep at least half of the annotated edges'
gain on the same notes. The learned weights, the other channels, the random
scheme and the pitch term were recorded without prediction.

\section{Development results}\label{sec:motif-dev}

\subsection{The shared form}\label{sec:motif-pilot1}
With whole occurrences hidden, joining their notes to the observed returns
of the same motif lowers the timing log score on those notes by $0.117$
nats against the time-only graph, interval $[-0.245, -0.006]$, on
twenty-one movements and two seeds. The pitch term of the thesis's graph
gives $-0.014$ on the same notes, interval $[-0.046, +0.016]$. The timing
error falls by $0.0026$~s, interval $[-0.0054, -0.0003]$, which is the gain
the pitch term also gives, so on timing accuracy the motif edges do what
the pitch term does and the extra is in calibration. The clearest gain is
articulation accuracy, root mean squared error $0.637$ to $0.611$, a
difference of $-0.033$ with interval $[-0.066, -0.007]$, with the log score
moving the same way but with an interval covering zero. Velocity gets
worse in calibration, log score $+0.069$, interval $[+0.005, +0.138]$, and
the coverage of the $90\%$ band falls from $0.878$ to $0.823$: the motif
edges tie the velocities of returns together more tightly than performers
play them, and one weight shared by the three channels cannot loosen one
channel alone. Under the random scheme the family is mostly invisible, as
expected. The evidence treats the family as a switch per piece, as it
treats the feature kernels: the learned weight falls below $0.1$ on eight
movements and exceeds $2$ on six.

\subsection{The per-channel form}\label{sec:motif-pilot2}
The second study repeated the first on the same movements and masks with
four seeds, the first two of which reproduce the first study exactly, and
added the per-channel form of \eqref{eq:motif-pc}. Its prediction was that
the per-channel form keeps the timing gain of the shared form and removes
the velocity miscalibration. The second half holds: the velocity log score
on the hidden motif notes falls by $0.086$ against the time-only graph,
interval $[-0.200, -0.018]$, the velocity error by $0.0057$, interval
$[-0.0106, -0.0023]$, and the coverage of the $90\%$ band rises from
$0.878$ to $0.944$. The first half does not: the per-channel form changes
the timing log score by $-0.005$, interval $[-0.076, +0.054]$. By the rule
fixed before the run, whatever timing value the shared form had came from
the coupling of the three channels through one weight, not from the edges
on their own.

The four seeds also change what the shared form's timing gain is. On all
four it is $-0.071$ with interval $[-0.230, +0.097]$, and on the two added
seeds alone $-0.026$ with interval $[-0.247, +0.259]$. The per-movement
differences have a median near zero, with three movements far below and one
far above, so the timing effect of the first study rested on two seeds and
three movements. The effect that holds on every pair of seeds and in both
forms is articulation accuracy: the root mean squared error on the hidden
motif notes falls from $0.732$ to $0.698$ in the shared form, a difference
of $-0.035$ with interval $[-0.060, -0.013]$, and to $0.713$ in the
per-channel form, $-0.019$ with interval $[-0.038, -0.005]$, with fifteen
and thirteen of the twenty-one movements below zero. The learned scales say
the same thing directly. On the timing channel $c_{\mathrm m}$ falls below
$0.01$ on nineteen movements, so the evidence switches the family off
where the time-only graph is already well calibrated. On the articulation
channel it exceeds $0.1$ on ten movements, and on the velocity channel it
is kept on seven. The two forms trade. The shared form gives the larger
articulation gain and a small timing error gain, $-0.0027$~s with interval
$[-0.0055, -0.0004]$, at the cost of velocity coverage, $0.835$. The
per-channel form gives a smaller articulation gain, velocity on both axes
and coverage at nominal, and nothing on timing. Adding the pitch term to
the per-channel form changes nothing. Under the random scheme the family
again has little to act on, and the per-channel form's timing log score
over all hidden notes is $+0.087$ with interval $[-0.000, +0.241]$, at the
edge of worse. The measured value of the family at this stage is therefore
articulation accuracy, and the pre-specified evaluation of
\S\ref{sec:motif-heldout} takes that as its primary measure.

\subsection{Predicted occurrences}\label{sec:motif-stage2}
The third study replaces the annotated occurrences by those predicted by
a graph-based motif discovery model \cite{lalang2026motif}, taking for each
movement the predictions of the cross-validation fold that held it out.
A predicted occurrence is a set of notes, mapped to the excerpt by onset
and pitch as the annotations were, with its span the first and last
matched onset and its type the model's pattern index. The masks and the
scored notes remain those of the annotated occurrences, identical to the
second study, so the edge source is the only change. At the note level the
predicted membership has a precision of $0.12$ to $1.00$ and a recall of
$0.44$ to $0.99$ against the annotation, with a median F1 of $0.78$. On two
movements the model predicts several times the annotated number of motif
notes and the precision falls below $0.2$. The prediction fixed before the
run was that the predicted edges keep at least half of the annotated
edges' timing log-score gain on the same notes, with an interval excluding
zero.

With the shared form the predicted edges lower the timing log score by
$0.034$, interval $[-0.085, +0.003]$, which is $0.48$ of the annotated
$0.071$ on the same notes with an interval that reaches zero, so the
prediction is not met, by a small margin on both counts. On articulation
accuracy, the effect that is robust in the second study, the predicted
edges give $-0.0135$ with interval $[-0.0317, +0.0030]$, four tenths of
the annotated $-0.0345$. The two edge sources act on the same movements:
their per-movement timing differences correlate at $0.59$ and their
articulation differences at $0.86$, and fifteen of the twenty-one
movements lie below zero for both. The movements with the lowest
membership precision receive large learned weights without a gain on the
annotated notes, while among the rest the size of the effect does not
order by F1. The per-channel form with predicted edges changes nothing and
costs $0.0014$~s of timing error, interval $[+0.0005, +0.0024]$. The
family can therefore be built from a model's occurrences on a score
without annotations and keeps about half of its effect at this size, with
intervals that cover zero on twenty-one movements.

\section{Held-out performances}\label{sec:motif-heldout}
% PENDING: the pre-specified evaluation (results/motif_heldout_results.md).

\section{What the result means}\label{sec:motif-meaning}
% PENDING: written after the held-out evaluation.

% =====================================================================
